\documentclass[11pt, a4paper]{article}

\usepackage[T1]{fontenc}
\usepackage[utf8]{inputenc}
\usepackage[english]{babel}
\usepackage{csquotes}
\usepackage[protrusion=true, expansion=false]{microtype}
\usepackage[left=1.3in, right=1.3in, top=1.1in, bottom=1.2in]{geometry}
\usepackage{xcolor}
\usepackage{parskip}
\usepackage{hyperref}
\usepackage{setspace}

% Colours
\definecolor{accent}{RGB}{130, 20, 35}       % dark burgundy red
\definecolor{rulegray}{RGB}{160, 160, 160}
\definecolor{taggray}{RGB}{100, 100, 100}

% No paragraph indent; comfortable inter-paragraph spacing
\setlength{\parindent}{0pt}
\setlength{\parskip}{0.75em}

% Stretch line spacing slightly for readability
\setstretch{1.12}

% Hyperref metadata
\hypersetup{
  colorlinks = true,
  urlcolor   = accent,
  linkcolor  = accent,
  pdftitle   = {Source Annotations — Andrei Dan Olaru},
  pdfauthor  = {Andrei Dan Olaru},
  pdfsubject = {MA Thesis Proposal, Uppsala University 2026},
}

% -----------------------------------------------------------------------
% Entry heading command:  \entry{number}{citation text}{tag}
% -----------------------------------------------------------------------
\newcommand{\entry}[3]{%
  \vspace{1.8em}
  {\color{rulegray}\rule{\linewidth}{0.4pt}}
  \vspace{0.5em}

  \noindent
  {\large\bfseries\color{accent}#1.}\enspace
  {\normalsize #2}%
  \enspace{\small\scshape\color{taggray}[#3]}

  \vspace{0.35em}
  {\color{accent}\rule{\linewidth}{1.2pt}}
  \vspace{0.2em}
}

% -----------------------------------------------------------------------
\begin{document}

% -----------------------------------------------------------------------
% Title block
% -----------------------------------------------------------------------
\begin{center}
  {\LARGE\bfseries\color{accent} Annotated Bibliography}\\[0.55em]
  {\large Andrei Dan Olaru}\\[0.25em]
  {\normalsize MA Thesis Proposal \textperiodcentered{} Uppsala University \textperiodcentered{} 2026}\\[0.2em]
  {\normalsize\itshape ``Applying Transformative Game Design Theory to Competitive Debate''}
\end{center}

{\color{accent}\rule{\linewidth}{1.8pt}}

% -----------------------------------------------------------------------
% Entry 1
% -----------------------------------------------------------------------
\entry{1}{Hugaas, K.\,H. (2024). Bleed and identity: A conceptual model of bleed and
how bleed-out from role-playing games can affect a player's sense of self.
\textit{International Journal of Role-Playing}, \textit{15}, 9--35.}{Theoretical Framework}

This journal article develops the most comprehensive conceptual model of bleed
available in role-playing game scholarship. The paper defines bleed as the
bidirectional spillover of mental and physical states, values, opinions, and
similar personal material between player and character, a term coined by
Emily Care Boss (2007) and refined here into a full theoretical framework. A central contribution is the bleed perception threshold: Hugaas
argues that bleed occurs in all players, but whether the player notices it
depends on whether it crosses their individual threshold, a distinction that
decouples bleed's transformative effects from the player's conscious awareness
of them. The paper distinguishes bleed-in (player to character) from bleed-out
(character to player) and focuses particularly on bleed-out's effects on the
player's sense of self. It synthesizes and extends earlier work into a taxonomy
of types including emotional, procedural, memetic, ego, identity, emancipatory,
and romantic bleed; identity bleed, where deeper self-schemas shift through
character inhabitation, is one of Hugaas's original contributions.

Bleed is the mechanism through which debate's documented transformative effects
are produced and explained in my thesis. The three types deployed most directly
are memetic bleed, the transfer of ideas and argumentative perspectives into the
debater's worldview outside the round; procedural bleed, the transfer of
rhetorical habits and communication skills into ordinary life; and emancipatory
bleed, where students from marginalized backgrounds develop intellectual voice
through the identity work of being taken seriously as arguers. The bleed
perception threshold concept is useful for the thesis's empirical grounding:
it explains why debaters often cannot articulate how debate changed them even
as longitudinal studies document consistent attitudinal shifts. For how debate
is broken, existing formats produce bleed without designing for it, leaving its
direction and content to competitive pressures rather than transformative goals.
For why the Inside Out format should work, its design addresses bleed at the
level of each type: the character assignment shapes memetic bleed, the debrief
converts bleed-out into conscious integration, and explicit identity-distance
mechanisms protect emancipatory bleed from being overridden by competitive
anxiety.

% -----------------------------------------------------------------------
% Entry 2
% -----------------------------------------------------------------------
\entry{2}{Daniau, S. (2016). The transformative potential of role-playing games:
From play skills to human skills. \textit{Simulation \& Gaming}, \textit{47}(4), 423--444.}{Theoretical Framework}

Drawing on a doctoral dissertation and ongoing action research, this
peer-reviewed SAGE journal article develops the concept of the Transformative
Role-Playing Game (TF-RPG): a tabletop RPG combined with a formal structured
debrief session facilitated after play. Daniau's core theoretical contribution
is that participants in RPGs simultaneously occupy four levels of reality
(character, player, person, and human being) and that the transfer of in-game
skills to lasting human competencies does not happen automatically through play
alone. It requires structured reflective processing during which participants
consciously map game experiences onto their own personal journeys. He builds a
four-dimensional learning framework (knowing, doing, being, and relating)
and argues that TF-RPGs uniquely engage all four, making them far more
educationally powerful than conventional instruction. Without the debrief, play
skills remain sealed inside the game. Daniau's four levels map the psychological
terrain of what happens inside the transformational container described by
Bowman et al.\ (2025): the container sets the conditions, but Daniau's debrief
is the mechanism that converts bleed-out into lasting change. His framework also
resonates with Kolb's (1984) experiential learning cycle, in which reflective
observation and abstract conceptualization (the stages where raw experience
becomes generalized learning) can only be activated by structured processing.

For my claim that debate is a role-playing game, Daniau's taxonomy of
role-playing types, which includes educational role-playing and simulation
alongside more traditional RPG formats, positions debate naturally within a
continuum of transformative role-playing activities. For how debate is broken,
Daniau's insight that the debrief is what converts play skills into human skills
becomes a pointed critique of current debate formats: the oral critique
processes argumentative performance (the character and player levels) but
entirely neglects the person and human being levels: there is no emotional
processing, no de-roling, and no structured connection between in-round
experience and the debater's developing self-concept. For why the Inside Out
format should work, I design an explicit debrief structure informed by Daniau
that addresses all four levels: it processes what arguments were made
(character), how the debater performed (player), what it felt like to argue
from that position (person), and what it means for how the debater understands
disagreement and persuasion in life (human being).

% -----------------------------------------------------------------------
% Entry 3
% -----------------------------------------------------------------------
\entry{3}{Bowman, S.\,L. (2010). \textit{The functions of role-playing games: How participants
create community, solve problems, and explore identity.} McFarland.}{Theoretical Framework}

This scholarly ethnography, based on participant-observer research and
extensive qualitative interviews with nineteen role-players across tabletop,
LARP, and virtual formats, identifies three core functions of role-playing:
communal cohesion (role-playing creates group identity through shared narrative
ritual), scenario-based problem solving (role-playing develops cognitive,
interpersonal, cultural, and professional skills through structured enactment
of complex situations), and identity alteration (role-playing provides a safe
space to enact alternate personas). The book's most theoretically generative
contribution for my purposes is Chapter~7's taxonomy of nine types of
character-player relationships, among them the Fragmented Self, the Oppositional
Self, the Experimental Self, and the Repressed Self, each describing a
different way that a player's own psychological content is activated, magnified,
or reorganized through the character. Central to all nine types is the concept
of adopting a theory of mind: thinking as if one were someone else in a unique
set of circumstances, which Bowman argues develops empathy and expands
participants' understanding of perspectives beyond their own. This framework
provides the foundational identity architecture that underpins the bleed
taxonomy in Bowman et al.\ (2025): bleed can only flow between player and
character because the character is built from the player's own psychological
material: fragments, repressions, ideals, experiments. Without the nine-type
identity structure, bleed would be a metaphor without a mechanism.

For my claim that debate is a role-playing game, Bowman's framework is
decisive: the debater's relationship to their assigned position is structurally
identical to what she calls the Fragmented Self relationship. The debater does
not pretend to believe the motion; they locate within themselves the fragments
of belief, knowledge, and rhetorical disposition that can genuinely support it,
then magnify those fragments into a complete epistemic stance. This is authentic
identity work, not theatrical performance, and this is why debate produces
real cognitive and attitudinal effects rather than merely simulated ones.
For how debate is broken, the communal cohesion function reveals a
double-edged dynamic: the tight sub-culture of competitive debate creates
strong in-group identity, but when that identity becomes hierarchically
structured around tournament success, the community begins to suppress rather
than expand its members' identities, blocking the Experimental Self and
Repressed Self relationships that would produce the most growth. For why the
Inside Out format should work, I design character assignments that deliberately
activate the Oppositional Self relationship: debaters are given not just a side
but a specific persona whose perspective differs meaningfully from their own,
and structured community practices strengthen the communal cohesion function
without generating competitive hierarchy.

% -----------------------------------------------------------------------
% Entry 4
% -----------------------------------------------------------------------
\entry{4}{Hammer, J., To, A., Schrier, K., Bowman, S.\,L., \& Kaufman, G. (2024).
Learning and role-playing games. In J.\,P. Zagal \& S. Deterding (Eds.),
\textit{The Routledge Handbook of Role-Playing Game Studies} (pp.~299--317).
Routledge.}{Theoretical Framework}

This book chapter provides a comprehensive synthesis of educational theory
as applied to role-playing games.
The authors situate RPGs within four learning theory traditions (behaviorism,
cognitivism, constructivism, and sociocultural theory) and argue that RPGs
operate primarily through constructionist and sociocultural mechanisms: they
are shared constructions of fictional reality requiring players to adopt and
inhabit new roles within a community of practice. The chapter's central
contribution is a taxonomy of five educational features distinctive to RPGs:
character portrayal (enabling perspective-taking, experience-taking, and
vicarious experience); manipulation of a fictional world through theorycrafting
and authentic simulation; experiencing an altered sense of reality, which can
either heighten the emotional stakes of fictional decisions or lower the stakes
of realistic ones; supporting collaboration through shared imaginative spaces;
and learning through making RPGs, including modding, designing new games, and
using game creation as critical making practice. Research synthesized in the
chapter shows that perspective-taking through character adoption increases
empathy and reduces reliance on stereotypes, while experience-taking (deeper
immersion where players simulate a character's subjective experience) produces
temporary internalization of the character's personality traits and behavior.
The chapter also identifies the debrief as the mechanism through which
in-game experiences are processed for learning, distinguishing emotional,
intellectual, and educational processing, and emphasizing that without
structured reflection, learning remains locked inside the game experience.

This chapter is the educational theory bridge between RPG scholarship and the
thesis's claims about debate. The five educational features of RPGs map onto
the debate encounter: switch-side argumentation is a form of character
portrayal that cultivates perspective-taking comparable to role adoption in
RPGs; the adversarial dialogue is a form of authentic simulation where
debaters experiment with how argumentative positions hold up under pressure;
team preparation is a form of collaborative shared imagination. Hammer et al.'s
evidence that immersion depth correlates with internalization of
character-perspective traits provides theoretical warrant for designing debate
formats that deepen character identification rather than treating position
assignment as a strategic convenience. Their debrief taxonomy clarifies why
the oral critique falls short: because educational processing requires
structured reflection connecting in-game experience to real-world
understanding, a critique focused exclusively on argumentative performance
does not activate the learning mechanisms the debate format has already
engaged.

% -----------------------------------------------------------------------
% Entry 5
% -----------------------------------------------------------------------
\entry{5}{Deterding, S. (2017). Alibis for adult play: A Goffmanian account of
escaping embarrassment in adult play. \textit{Games and Culture}, \textit{13}(3), 260--279.}{Theoretical Framework}

This journal article develops the concept of alibi from Erving Goffman's
sociology of face-work and frame analysis. In Goffman's account, face is the
positive social value a person claims for themselves in interaction; adult life
is governed by norms of seriousness that make play a potential threat to face,
liable to make the player appear childish or frivolous. Deterding argues that
adult play persists because participants deploy alibis: motivational accounts
that deflect negative inference from displayed behavior to a person's identity.
The paper identifies two sources of alibi. First, motivational accounting
systems, in which players reframe play through adult-appropriate purposes such
as professional development, family engagement, or health: play becomes
legitimate because it serves a serious end. Second, keyings or reframings, in
which play is bracketed as non-serious through ironic, humorous, or
self-distancing performance that signals the player's distance from the
activity (deliberately bad karaoke, for instance, preemptively frames
mediocre performance as intentional). Together these mechanisms create
plausible deniability: the reason for engaging in the behavior is not
\enquote{I want to play} but \enquote{I have good adult reasons for this activity} or
\enquote{I am not really doing this seriously.}

In the thesis, alibi is the concept that explains why debate's epistemic
character inhabitation is psychologically safe: debaters can argue for positions
contrary to their own beliefs because the competitive frame provides alibi. The
random side assignment is the mechanism: the reason a debater argues for a
position is \enquote{I was assigned this side,} not \enquote{I personally believe this,}
which deflects inference from displayed behavior to the game's structure. This
is what Seo (2022) describes as debate providing a \enquote{sense of freedom} to
\enquote{flirt with ideas, unencumbered by expectations of consistency or deep
conviction.} For how debate is broken, when competitive pressures make
tournament success feel like a measure of personal worth, the alibi erodes:
participants can no longer treat debate as just a game because the stakes have
become real, and the protected space for identity risk-taking collapses. For
why the Inside Out format should work, it calibrates alibi deliberately: enough
protective bracketing for participants to take genuine identity risks, but with
debrief practices that thin the alibi at specific points so bleed-out can enter
conscious reflection rather than dissipate.

% -----------------------------------------------------------------------
% Closing rule
% -----------------------------------------------------------------------
\vspace{1.8em}
{\color{rulegray}\rule{\linewidth}{0.4pt}}

\end{document}
